\documentclass[12pt,a4paper]{article}
\usepackage[utf8]{inputenc}
\usepackage[T1]{fontenc}
\usepackage[spanish]{babel}
\usepackage{amsmath,amssymb}
\usepackage{geometry}
\usepackage{booktabs}
\usepackage{longtable}
\usepackage{graphicx}
\usepackage{array}
\usepackage{pdflscape}
\geometry{margin=2.5cm}

\title{Tarea --- Análisis Comparativo de Conductores Eléctricos\\para Líneas de Transmisión}
\author{Carla\\Líneas de Transmisión}
\date{1 de junio de 2026}

\begin{document}
\maketitle

\section{Introducción}

Este informe presenta un análisis comparativo de los principales tipos de
conductores utilizados en líneas de transmisión eléctrica, evaluando sus
parámetros técnicos para el calibre asignado de \textbf{1100 MCM}.

La investigación se realizó con base en el catálogo técnico de \textbf{CVG CABELUM}
(Conductores de Aluminio del Caroní, C.A.), empresa estatal venezolana ubicada
en Ciudad Bolívar, única productora nacional de conductores eléctricos de alta
potencia. La sección transversal aproximada para 1100 MCM es de
$557.37\text{ mm}^2$.

\section{Consideraciones de Fabricación y Alcance del Catálogo}

Al revisar exhaustivamente el catálogo de CVG CABELUM, se identificaron las
siguientes consideraciones técnicas importantes:

\begin{enumerate}
    \item \textbf{Calibres exactos:} El catálogo fabrica de forma exacta el
        calibre \textbf{1100.0 MCM} para las tecnologías \textbf{AAAC} (pág.
        32) y \textbf{ACAR} (pág. 44). Para \textbf{AAC} y \textbf{ACSR} los
        calibres comerciales más cercanos son \textbf{1113 MCM} (\textit{Marigold}
        en AAC; \textit{Finch / Bluejay} en ACSR).

    \item \textbf{Normas DIN y BRITISH:} Estas normas no miden en MCM, sino
        en mm$^2$. Para el valor de 1100 MCM ($\approx 557\text{ mm}^2$), la
        planta tiene dos calibres comerciales: \textbf{DIN 550/560}
        ($549.65\text{--}561.70\text{ mm}^2$) y el calibre BRITISH
        \textbf{Moose} ($500\text{ mm}^2$ nominal).

    \item \textbf{Limitaciones:} La norma BRITISH se queda corta en AAC y AAAC
        (topando en calibres menores de 750/700 mm$^2$), pero en ACSR dispone
        del calibre \textbf{Moose}. Las normas DIN y BRITISH no contemplan la
        tipología \textbf{ACAR}.
\end{enumerate}

\section{Tipos de Conductores Evaluados}

\subsection{AAC (All Aluminum Conductor) --- ASTM B-231}
Conductor de aluminio puro 1350-H19. Características:
\begin{itemize}
    \item Alta conductividad eléctrica
    \item Bajo peso
    \item Menor resistencia mecánica
    \item Usos: distribución urbana, cortas distancias
\end{itemize}

\subsection{AAAC (All Aluminum Alloy Conductor) --- ASTM B-399}
Conductor de aleación de aluminio 6201-T81. Características:
\begin{itemize}
    \item Mayor resistencia mecánica que AAC
    \item Buena relación resistencia-peso
    \item Resistencia a la corrosión
    \item Usos: líneas de media y larga distancia
\end{itemize}

\subsection{ACAR (Aluminum Conductor Alloy Reinforced) --- ASTM B-524}
Conductor de aluminio con refuerzo de aleación. Características:
\begin{itemize}
    \item Combinación de alambres de aluminio puro y de aleación 6201
    \item Balance entre conductividad y resistencia
    \item Usos: líneas que requieren mayor capacidad mecánica
\end{itemize}

\subsection{ACSR (Aluminum Conductor Steel Reinforced) --- ASTM B-232}
Conductor de aluminio con refuerzo de acero galvanizado. Características:
\begin{itemize}
    \item El más utilizado en transmisión
    \item Excelente relación resistencia-peso
    \item Máxima resistencia mecánica de la línea
    \item Usos: líneas de transmisión de alta tensión
\end{itemize}

\section{Tabla Comparativa --- Calibre 1100 MCM / 1113 MCM}

Evaluado bajo los parámetros disponibles en el catálogo de CABELUM para el
calibre asignado:

\begin{landscape}
\begin{center}
\begin{tabular}{p{4cm}c|c|c|c}
\toprule
\textbf{Parámetro} & \textbf{AAC} & \textbf{AAAC} & \textbf{ACAR} & \textbf{ACSR} \\
& (Marigold) & (1100 MCM) & (1100 MCM 33/4) & (Finch) \\
\midrule
Norma de fabricación & ASTM B-231 & ASTM B-399 & ASTM B-524 & ASTM B-232 y B-498 \\
\midrule
Calibre comercial & 1113 kcmil & 1100 kcmil & 1100 kcmil & 1113 kcmil \\
\midrule
Sección (mm$^2$) & 564.0 & 557.0 & 557.0 & 636.60 \\
\midrule
Diámetro (mm) & 30.87 & 30.42 & 30.65 & 32.85 \\
\midrule
Peso neto (kg/km) & 1553.0 & 1502 & 1532 & 2128 \\
\midrule
Carga de rotura (kgf) & 8903 & 15909 & 9413 & 17737 \\
\midrule
Resistencia DC 20$^\circ$C ($\Omega$/km) & 0.05112 & 0.06120 & 0.05264 & 0.05152 \\
\midrule
Ampacidad (A) & 1080 & 998 & 1061 & 1100 \\
\midrule
Long. máx. despacho (m) & 2120 & 2770 & 2150 & 1580 \\
\midrule
Peso neto carrete (kg) & 3292 & 4141 & 3294 & 3362 \\
\midrule
Peso bruto carrete (kg) & 3592 & 4571 & 3594 & 3662 \\
\midrule
Tipo de carrete & $75"\times38"\times34"$ & $84"$ & $75"\times38"\times34"$ & $75"\times38"\times34"$ \\
\midrule
Marca & CABELUM & CABELUM & CABELUM & CABELUM \\
\bottomrule
\end{tabular}
\end{center}

\vspace{0.3cm}
\textbf{Nota sobre ampacidad:} Calculada bajo condiciones estándar de CABELUM
(temperatura ambiente 25$^\circ$C, temperatura del conductor 75$^\circ$C,
viento de 0.6 m/s y radiación solar de 1000 W/m$^2$).
\end{landscape}

\section{Desglose por Norma y Tipo de Conductor}

\begin{landscape}
\subsection{Conductor AAC (All Aluminum Conductor)}

Comparativa del calibre Marigold (ASTM) frente a los calibres más cercanos
de las normas europeas:

\begin{center}
\begin{tabular}{p{4cm}c|c|c}
\toprule
\textbf{Parámetro} & \textbf{ASTM B-231} & \textbf{BRITISH (BS 215)} & \textbf{DIN 48201-5} \\
& (1113 kcmil Marigold) & (Tarantula) & (500 mm$^2$) \\
\midrule
Designación de calibre & 1113 kcmil & 750 mm$^2$ nominal & 500 mm$^2$ nominal \\
\midrule
Sección real (mm$^2$) & 564.0 & 794.87 & 499.83 \\
\midrule
Diámetro (mm) & 30.87 & 36.61 & 29.10 \\
\midrule
Peso (kg/km) & 1553.0 & 2192 & 1379 \\
\midrule
Carga de rotura (kgf) & 8903 & 12243 & 7615 \\
\midrule
Resistencia DC 20$^\circ$C ($\Omega$/km) & 0.05112 & 0.03627 & 0.05782 \\
\midrule
Ampacidad (A) & 1080 & 899 & 689 \\
\midrule
Long. máx. despacho (m) & 2120 & 1510 & 1120 \\
\midrule
Peso neto carrete (kg) & 3292 & 3310 & 1544 \\
\midrule
Peso bruto carrete (kg) & 3592 & 3610 & 1689 \\
\midrule
Marca & CABELUM & CABELUM & CABELUM \\
\bottomrule
\end{tabular}
\end{center}
\end{landscape}

\begin{landscape}
\subsection{Conductor AAAC (All Aluminum Alloy Conductor)}

Comparativa del calibre exacto 1100 MCM (ASTM) frente a los límites europeos:

\begin{center}
\begin{tabular}{p{4cm}c|c|c}
\toprule
\textbf{Parámetro} & \textbf{ASTM B-399} & \textbf{BRITISH (BS 3242)} & \textbf{DIN 48201-6} \\
& (1100 kcmil) & (Araucaria) & (500 mm$^2$) \\
\midrule
Designación de calibre & 1100 kcmil & 700 mm$^2$ nominal & 500 mm$^2$ nominal \\
\midrule
Sección real (mm$^2$) & 557.0 & 821.15 & 499.83 \\
\midrule
Diámetro (mm) & 30.42 & 37.26 & 29.10 \\
\midrule
Peso (kg/km) & 1502 & 2266 & 1379 \\
\midrule
Carga de rotura (kgf) & 15909 & 23465 & 14236 \\
\midrule
Resistencia DC 20$^\circ$C ($\Omega$/km) & 0.06120 & 0.04046 & 0.0671 \\
\midrule
Ampacidad (A) & 998 & --- & --- \\
\midrule
Long. máx. despacho (m) & 2770 & 1450 & 2390 \\
\midrule
Peso neto carrete (kg) & 4141 & 3286 & 3296 \\
\midrule
Peso bruto carrete (kg) & 4571 & 3586 & 3596 \\
\midrule
Marca & CABELUM & CABELUM & CABELUM \\
\bottomrule
\end{tabular}
\end{center}
\end{landscape}

\begin{landscape}
\subsection{Conductor ACAR (Aluminum Conductor Alloy Reinforced)}

La norma BRITISH y la norma DIN \textbf{no contemplan la fabricación de la
tipología ACAR} en el catálogo de CABELUM. Se fabrica bajo pedido únicamente
con estándares americanos.

\begin{center}
\begin{tabular}{p{4cm}c|c|c}
\toprule
\textbf{Parámetro} & \textbf{ASTM B-524} & \textbf{BRITISH} & \textbf{DIN} \\
& (1100.0 MCM 33/4) & & \\
\midrule
Designación de calibre & 1100.0 MCM (33 Al / 4 Aleac.) & No disponible & No disponible \\
\midrule
Sección real (mm$^2$) & 557.0 & --- & --- \\
\midrule
Diámetro (mm) & 30.65 & --- & --- \\
\midrule
Peso (kg/km) & 1532 & --- & --- \\
\midrule
Carga de rotura (kgf) & 9413 & --- & --- \\
\midrule
Resistencia DC 20$^\circ$C ($\Omega$/km) & 0.05264 & --- & --- \\
\midrule
Ampacidad (A) & 1061 & --- & --- \\
\midrule
Long. máx. despacho (m) & 2150 & --- & --- \\
\midrule
Peso neto carrete (kg) & 3294 & --- & --- \\
\midrule
Peso bruto carrete (kg) & 3594 & --- & --- \\
\midrule
Marca & CABELUM & --- & --- \\
\bottomrule
\end{tabular}
\end{center}
\end{landscape}

\begin{landscape}
\subsection{Conductor ACSR (Aluminum Conductor Steel Reinforced)}

Comparativa del modelo Finch (1113 kcmil) contra los calibres de gran
potencia europeos:

\begin{center}
\begin{tabular}{p{4cm}c|c|c}
\toprule
\textbf{Parámetro} & \textbf{ASTM B-232} & \textbf{BRITISH (BS 215 P.2)} & \textbf{DIN 48204} \\
& (1113 kcmil Finch) & (Moose) & (550/70) \\
\midrule
Designación de calibre & 1113 kcmil & 500 mm$^2$ nominal & 550/70 mm$^2$ nominal \\
\midrule
Sección real (mm$^2$) & 636.60 & 597.00 & 620.90 \\
\midrule
Diámetro (mm) & 32.85 & 31.77 & 32.40 \\
\midrule
Peso (kg/km) & 2128 & 1996 & 2085 \\
\midrule
Carga de rotura (kgf) & 17737 & 16417 & 17073 \\
\midrule
Resistencia DC 20$^\circ$C ($\Omega$/km) & 0.05152 & 0.05470 & 0.0526 \\
\midrule
Ampacidad (A) & 1100 & 763 & 734 \\
\midrule
Long. máx. despacho (m) & 1580 & 1690 & 1620 \\
\midrule
Peso neto carrete (kg) & 3362 & 3377 & 3378 \\
\midrule
Peso bruto carrete (kg) & 3662 & 3677 & 3678 \\
\midrule
Marca & CABELUM & CABELUM & CABELUM \\
\bottomrule
\end{tabular}
\end{center}
\end{landscape}

\section{Análisis Comparativo}

\subsection{Por tipo de conductor}

El \textbf{AAC Marigold} (1113 kcmil) presenta la menor resistencia eléctrica
(0.05112 $\Omega$/km) y una ampacidad de 1080 A, pero su carga de rotura
(8903 kgf) es la más baja, limitando su uso a cortas distancias.

El \textbf{AAAC 1100 MCM} ofrece la carga de rotura más alta del grupo
(15909 kgf) gracias a la aleación 6201-T81, aunque su resistencia eléctrica
es la mayor (0.06120 $\Omega$/km) y su ampacidad la más baja (998 A).

El \textbf{ACAR 1100 MCM 33/4} consigue una ampacidad de 1061 A con una
carga de rotura de 9413 kgf, siendo una opción equilibrada.

El \textbf{ACSR Finch} (1113 kcmil), con su núcleo de acero, ofrece la mayor
carga de rotura (17737 kgf) y la ampacidad más alta (1100 A), aunque su peso
(2128 kg/km) es significativamente mayor.

\subsection{Por norma de fabricación}

Las normas \textbf{DIN} y \textbf{BRITISH} no contemplan la denominación MCM.
Para el calibre de Carla (1100 MCM $\approx 557\text{ mm}^2$), el calibre
europeo más cercano es el \textbf{DIN 500 mm$^2$} en AAC y AAAC, y el
\textbf{Moose 500 mm$^2$} (BS) o \textbf{550/70 mm$^2$} (DIN) en ACSR.

\subsection{Conductor asignado: Carla 1100 MCM}

El conductor de tipo \textbf{AAAC 1100 MCM} ofrece la mayor resistencia
mecánica del grupo (15909 kgf), con una sección de 557 mm$^2$ y un diámetro
de 30.42 mm, fabricado por CVG CABELUM bajo norma ASTM B-399.

\section{Contacto con Proveedor --- CVG CABELUM}

\subsection{Datos de Contacto}

\begin{center}
\begin{tabular}{ll}
\toprule
\textbf{Dependencia} & \textbf{Datos} \\
\midrule
Oficina Comercial Caracas & Av. La Estancia, Edf. General, Piso 4, Chuao \\
& Teléfonos: 0414-103.16.55 / 0416-834.25.29 \\
\midrule
Planta Ciudad Bolívar & Vía Perimetral, Nueva Zona Industrial \\
& Teléfonos: 0285-444.35.22 / 0416-685.12.18 \\
\midrule
Correos electrónicos & comercializacion@cabelum.com.ve \\
& contacto@cabelum.com.ve \\
\bottomrule
\end{tabular}
\end{center}

\subsection{Condiciones Comerciales}

\begin{center}
\begin{tabular}{ll}
\toprule
\textbf{Aspecto} & \textbf{Información} \\
\midrule
Precio & No es fijo. Se cotiza bajo Fórmula de Reajuste \\
& de Precios indexada al valor del metal en la \\
& bolsa de la LME (London Metal Exchange). \\
\midrule
Condiciones de pago & Anticipo del 50\% al 70\% para proceder al \\
& trefilado y cableado. Saldo restante contra \\
& informes de inspección de calidad en fábrica. \\
\midrule
Condiciones de entrega & FOB Planta (Incoterms) en Ciudad Bolívar u \\
& O.A.D (Puesto en Obra) según se acuerde. \\
\midrule
Plazo de entrega & 45 a 60 días, dependiendo de la cola de \\
& producción y disponibilidad de alambrón. \\
\bottomrule
\end{tabular}
\end{center}

\section{Conclusiones Técnicas}

\begin{enumerate}
    \item \textbf{Comportamiento eléctrico frente al mecánico:} El conductor de
        tipo \textbf{AAAC 1100 MCM} ofrece la mayor carga de rotura del grupo
        (15909 kgf), siendo la opción más robusta mecánicamente, aunque su
        ampacidad (998 A) es la menor del conjunto.

    \item \textbf{Calibres exactos:} CABELUM fabrica el calibre exacto
        \textbf{1100 MCM} en AAAC y ACAR. Para AAC y ACSR los calibres
        comerciales más cercanos son \textbf{1113 MCM} (Marigold y Finch).

    \item \textbf{Limitaciones de normas:} Las normas \textbf{DIN} y
        \textbf{BRITISH} no contemplan la denominación MCM. El calibre DIN
        más cercano es el 500 mm$^2$ y el BRITISH el Moose (500 mm$^2$).

    \item \textbf{Logística de empaque:} Los carretes utilizados son modelos
        de 75" × 38" × 34" o de 84" para el AAAC (4500 Kg), limitando la
        estiba a 6 o 4 carretes por contenedor de 20 pies respectivamente.
\end{enumerate}

\section{Referencias}

\begin{enumerate}
    \item ASTM B-231 --- Especificación estándar para conductores AAC de aluminio 1350
    \item ASTM B-399 --- Especificación estándar para conductores AAAC de aleación 6201-T81
    \item ASTM B-524 --- Especificación estándar para conductores ACAR
    \item ASTM B-232 --- Especificación estándar para conductores ACSR
    \item CVG CABELUM --- Catálogo técnico de conductores desnudos, Planta Ciudad Bolívar
    \item British Standards BS 215 / BS 3242 --- Conductores para líneas de transmisión
    \item DIN 48201 / 48204 --- Normas alemanas para conductores eléctricos
\end{enumerate}

\end{document}
